\section{Introduction}
A persistent-memory agent can diverge from itself twice: first when feedback changes what it stores from an otherwise identical experience, and again when only some of that stored difference survives into later behavior. Consider two copies of the same interaction trajectory. If one copy is routed through the released success-reflection branch and the other through the failure-reflection branch, the writer can store different reusable advice even though the underlying experience is unchanged. The first scientific question is therefore \emph{how strongly does feedback-conditioned writing change persistent state?} The second is \emph{once that difference exists, where does the deployed reuse chain preserve or lose it before behavior changes?}

These two questions belong in the same paper because persistent self-improvement contains several causal surfaces that are often collapsed into one ``memory helped'' statement. A terminal signal may first alter what is \textbf{written}; the resulting branch-conditioned source-memory item may or may not be \textbf{exposed} by retrieval on a later task; the policy may or may not \textbf{use} the treatment-specific information within that retrieved item; and only then can behavior change at the terminal \textbf{outcome}. A positive memory-similarity result identifies the first surface, while a retrieval hit only identifies availability. Neither is evidence by itself that the agent's decision rule changed. Conversely, an endpoint null does not reveal whether the update failed to write, failed to retrieve, or was filtered after retrieval. We call this \emph{evaluation coarsening}: observing only a projection of the stage chain aliases distinct internal explanations into the same coarse success or failure label.

We study this problem through a stage-resolved intervention rather than a new memory algorithm. For a frozen trajectory $\tau$, we switch the released success/failure reflection branch and construct paired persistent memories while holding the source experience fixed. This is a bundled writer-protocol intervention: the branch changes the reflection instruction together with its outcome/reward semantics, so it is not an atom-level isolation of a reward bit. We then instrument four native stages,
\[
\text{write} \rightarrow \text{retrieval exposure} \rightarrow \text{first-action uptake} \rightarrow \text{terminal outcome},
\]
and add a separate forced fixed-evidence control that bypasses native retrieval. The forced arm is deliberately \emph{not} another stage in the chain: it asks whether branch-specific memory content has downstream leverage when supplied, whereas the native path asks whether the deployed memory system actually transports that difference into behavior.

This design distinguishes several explanations that an endpoint-only experiment cannot separate. If the writer does not create a persistent intervention, the story stops at write time. If the downstream policy is globally insensitive to memory content, even forced exposure should have little leverage. If the only problem is retrieval absence, native exposure should collapse. If retrieval itself is a valid proxy for policy use, observed exposure should co-occur with a supported action-level branch contrast. Finally, if reward-conditioned memory has a universal directional effect, terminal differences should not be dominated by zeros and sign reversals across tasks. Because these stages have different units, sample units, tests, and experimental surfaces, we do not convert them into a single ``transport efficiency.'' Instead, we retain a \emph{stage signature} and introduce an ordinal stage-evidence ladder: preserve each stage's own inferential/descriptive semantics, then localize the first \emph{measured} native stage whose branch contrast is no longer supported after the preceding prerequisites have been supported or directly observed. The stage tests are not power-matched, so this operator localizes an evidence boundary rather than the latent point where a causal effect must physically decay. The point is diagnostic disambiguation rather than a finer-grained score.

The frozen evidence yields a sharp pattern. Reward-conditioned writing changes durable memory in all 20 Shopping pairs and all 4 Reddit pairs. A same-mode wording control shows that the reward/reflection-mode contrast exceeds a stronger within-mode wording perturbation. Forced fixed-evidence exposure produces terminal $|\Delta|=0.15625$ ($p=0.00074$), demonstrating downstream leverage under supplied memory. Yet native Shopping retrieval remains common---125 hits in 172 held-out opportunities---while first-action TV is only $0.06944$ ($p=0.5801$) with 0/36 modal action changes, and native terminal $|\Delta|$ is $0.02083$ ($p=0.4289$) with 34/36 zero cells. Reddit again shows reliable write divergence but sparse terminal transport: 6/8 cells are zero and the two nonzero effects have opposite signs.

Taken together, these measurements weaken three simple explanations at once: the native null is not caused by failure to create persistent state; it is not well described as global downstream insensitivity to supplied memory content; and it is not explained by retrieval absence alone. More importantly, they expose why coarse evaluations disagree. Write-only evaluation would call the update strong; retrieval-only evaluation would call reuse active; native endpoint-only evaluation would call the effect negligible; forced-only evaluation would call the memory behaviorally consequential. The full stage signature resolves this apparent contradiction: on the frozen Shopping tests, first-action uptake is the first unsupported measured native stage after durable write and source-item exposure have been established. We intentionally stop short of a causal mediation claim because the frozen pool lacks per-atom decoding-seed binding and same-condition noise-floor replication, and the forced and native surfaces are not identical interventions.

Our contributions are:
\begin{enumerate}
    \item \textbf{A controlled feedback-to-state result.} Holding the source trajectory fixed, we switch the released success/failure reflection branch as a bundled writer-protocol intervention and show robust persistent-state divergence across Shopping and Reddit, including a stronger same-mode wording control that does not explain the write effect away.
    \item \textbf{A stage-resolved account of what happens after writing.} We follow that same persistent-state contrast through native source-item exposure, first-action uptake, and terminal outcome rather than treating state divergence as behavioral authority.
    \item \textbf{A capacity-versus-transport control.} Forced fixed-evidence exposure establishes that branch-specific memory can have downstream leverage when supplied, while native reuse tests whether the deployed retrieval-and-policy path actually realizes that leverage.
    \item \textbf{A diagnostic stage signature and cross-domain boundary.} We show that write-only, retrieval-only, native-endpoint-only, and forced-only views support different coarse summaries, and localize the first unsupported measured native Shopping stage without inventing a cross-stage attenuation ratio or causal mediation coefficient. Reddit reproduces write-stage divergence but not a stable signed native terminal effect.
\end{enumerate}

The resulting claim is deliberately narrower than a new memory method: \textbf{persistent-state divergence is not behavioral authority, and retrieval availability is not policy use}. This boundary tells us what a self-improving agent evaluation must measure before crediting a persistent update with downstream control.
